\begin{trapbox}
	\begin{description}[style=nextline, leftmargin=2.8em, labelsep=0.5em, itemsep=0.35em]
		\item[\textbf{\textcolor{CTrap}{易错点一（参数大小混淆）}}] 误认为 $c$ 是最大边，或混淆 $a$、$b$、$c$ 的大小顺序。
		\item[\textbf{\textcolor{CTrap}{正确理解（a恒为最大）：}}] 在椭圆中，$a$ 是长半轴长，满足 $a^{2}=b^{2}+c^{2}$，因此 $a$ 最大，$c$ 可能大于 $b$ 也可能小于 $b$。
		\item[\textbf{\textcolor{CTrap}{错因分析（关系混淆）：}}] 容易将椭圆的关系与双曲线的 $c^{2}=a^{2}+b^{2}$ 混淆，导致大小关系错误。
	\end{description}

	\medskip
	\begin{description}[style=nextline, leftmargin=2.8em, labelsep=0.5em, itemsep=0.35em]
		\item[\textbf{\textcolor{CTrap}{易错点二（焦点位置误判）}}] 当椭圆方程为 $\frac{y^{2}}{a^{2}}+\frac{x^{2}}{b^{2}}=1$ 时，误认为 $x^{2}$ 分母是 $a^{2}$，从而取错 $a$、$b$。
		\item[\textbf{\textcolor{CTrap}{正确理解（看分母大小）：}}] 标准方程中，哪个分母大，哪个就是 $a^{2}$，与变量名称无关。
		\item[\textbf{\textcolor{CTrap}{错因分析（思维定式）：}}] 习惯认为 $x^{2}$ 下方为 $a^{2}$，忽略了焦点在 $y$ 轴的情形。
	\end{description}

	\medskip
	\begin{description}[style=nextline, leftmargin=2.8em, labelsep=0.5em, itemsep=0.35em]
		\item[\textbf{\textcolor{CTrap}{易错点三（开方符号忽略）}}] 由 $c^{2}=a^{2}-b^{2}$ 开方求 $c$ 时，写出 $c=\pm\sqrt{a^{2}-b^{2}}$，而椭圆中的 $c$ 仅为正值。
		\item[\textbf{\textcolor{CTrap}{正确理解（参数正数要求）：}} $a,b,c$ 均为正数，开方只取算术平方根，即 $c=\sqrt{a^{2}-b^{2}}$。
		\item[\textbf{\textcolor{CTrap}{错因分析（习惯性考虑负根）：}}] 在代数解方程时习惯保留 $\pm$，但几何参数规定了正数范围，需舍去负根。
	\end{description}
\end{trapbox}
